% #### FRONTESPIZIO (frontmatter) ####
\setlength{\headheight}{33pt} %Distanzia l'header
\pagenumbering{gobble} %Toglie il numero di pagina
\begin{titlepage}
	\begin{center}
		\includegraphics[scale=0.6]{img/logo.png} \\ %Logo
		\vspace{0.4cm} %Aggiunge uno spazio verticale di 0.5 cm
		
		{\LARGE Progetto "Predire in Grafana"} \\ %Nome progetto
		\vspace{0.4cm} %Attenzione a mettere il punto e NON la virgola
		
		{\Huge \textbf{\DocTitle}} \\ %Titolo, prende variabile definita in metadata.tex
		\vspace{0.4cm}
		
		\DocDate \\ %Data, prende variabile definita in metadata.tex
		\vspace{0.4cm}
		
		%Allineamento colonne: l=left r=right c=center, 
		%va specificato per ogni colonna
		%Se si vuole la riga tra colonne mettere "|"
		
		\begin{tabular}{r | l} %Elementi colonne separate da "&", le righe finiscono con "\\"
			Versione             & \ver \\
			Approvazione         & \app \\ 
			Redazione            & \red \\
			Verifica             & \test \\
			Stato                & \stat \\
			Uso                  & \use \\
		    Destinato a          & Zucchetti \\
						         & Prof. Tullio Vardanega \\
						         & Prof. Riccardo Cardin \\
			Email di riferimento & vram.software@gmail.com
		\end{tabular}
		\vfill
		\textbf{Descrizione} \\
		\DocDesc
	\end{center}
\end{titlepage}
\clearpage

% #### Impostazione header, footer  e numerazione pagine ####
\pagenumbering{arabic} %Pagine con i numeri arabi + reset a 1
\renewcommand{\footrulewidth}{0.4pt} %Di default footrulewidth==0 e quindi è invisibile, di default \headrulewith==0.4pt
\rfoot{\thepage\ di \pageref{LastPage}} %Pagina n di m, con numeri Arabi; usa il pacchetto "lastpage", in caso non sia possibile usare tale pacchetto mettere al fondo dell'ultima pagina "\label{LastPage}"

% #### Tabella dei log ####
% \textbf = grassetto; \Large = font più grande
% \rowcolors{quanti colori alternare}{colore numero riga pari}{colore numero riga dispari}: colori alternati per riga
% \rowcolor{color}: cambia colore di una riga
% p{larghezza colonna}: p è un tipo di colonna di testo verticalmente allineata sopra, ci sarebbe anche m che è centrata a metà ma non è precisa per questo utilizzo TBStrut; la sintassi >{\centering} indica che il contenuto della colonna dovrà essere centrato
% \TBstrut fa parte di alcuni comandi che ho inserito in config.tex che permetto di aggiungere un po' di padding al testo
% \\ [2mm] : questra scrittura indica che lo spazio dopo una break line deve essere di 2mm
% 

%\setcounter{secnumdepth}{0}
%\hfill \break
%\textbf{\Large{Diario delle modifiche}} \\


\addtocontents{toc}{\protect\setcounter{tocdepth}{0}} %Inserire questo per escludere una sezione dall'indice.

\section*{Registro delle modifiche} %Asterisco per fare sezione non numerata
\rowcolors{2}{gray!25}{gray!15}
\begin{longtable} {
		>{\centering}p{17mm} 
		>{\centering}p{19.5mm}
		>{\centering}p{24mm} 
		>{\centering}p{24mm} 
		>{}p{32mm}}
	\rowcolor{gray!50}
	\textbf{Versione} & \textbf{Data} & \textbf{Nominativo} & \textbf{Ruolo} & \textbf{Descrizione} \TBstrut \\
	14.7.0 & 2020-04-09 & Stantagiuliana Vittorio, Toffoletto Massimo e Spreafico Alessandro & \textit{Progettista}, \textit{Verificatore} e \textit{Responsabile di progetto} & Stesura, verifica e approvazione documento. \TBstrut \\ [2mm]
\end{longtable}

\addtocontents{toc}{\protect\setcounter{tocdepth}{4}} %Inserire questo per ripristinare il normale inserimento delle sezioni nell'indice. 4 significa fino al paragrah
\clearpage

% #### INDICE (tableofcontents) ####
\tableofcontents %Provoca la stampa dell'indice
\clearpage

\setcounter{secnumdepth}{4} %Permette di andare fino alla profondità del paragraph con la numerazione delle sezioni